\documentclass[11pt,a4paper]{article}

\input{../preambule}

\begin{document}

\title{Semaine 5}
\author{}
\date{02/11 - 08/11 }

\maketitle

\begin{exercise}
Trouver toutes les fonctions $f:\N\to\N$ telles que
$$\forall n\in\N, f(n)+f(f(n))+f(f(f(n)))=3n$$
\end{exercise}

\begin{exercise}
On souhaite trouver toutes les fonctions $f:\R\to\R$ surjectives et strictement croissantes
 telles que $\forall x\in\R, f(x)+f^{-1}(x)=2x$.
\begin{questions}
\item Montrer que $f^{-1}$ est bien définie et exhiber une solution particulière.
\item On suppose qu'il existe $t\in\R$ tel que $f(t)\gt t$. On pose $a=f(t)-t$.
    \begin{questions}
    \item Montrer que $\forall n\in\Z, f(t+na)=t+(n+1)a$.
    \item Soit $x\in\R$ tel que $x$ n'est pas de la forme $t+na, n\in\Z$.
        \begin{questions}
        \item Justifier que $f(x)\gt x$.
        \item On note $b=f(x)-x$. Montrer que $\forall n\in\Z, f(x+nb)=x+(n+1)b$.
        \item Montrer que $a=b$.
        \end{questions}
    \end{questions}
\item Conclure.
\end{questions} 
\end{exercise}

\begin{exercise}
Trouver toutes les relations d'équivalence $R$ sur $\N$ telles que: 
$$\forall n,m\in\N, nRm\implies (n+1)R(m+1)$$
\end{exercise}

\begin{exercise}
Soit $A,B$ des ensembles.
On définit la relation $\preceq$ sur les cardinaux des ensembles par:
$$|A|\preceq |B|\text{ ssi il existe une injection de $A$ dans $B$}$$
\begin{questions}
\item Montrer que $\preceq$ est une relation d'ordre.
\item Montrer que $|\N^{\N}|=|\R|=|P(\N)|$.
\item Que dire de $|\R|$ et $|\R^{\R}|$
\end{questions}
\end{exercise}

\begin{exercise}
Soit $f:\C^*\to \C$ définie par $\forall z\in\C^*, f(z)=\frac{1}{z}$.
\begin{questions}
\item $f$ est-elle injective? surjective?
\item On considère deux ensembles:
$$E=\set{z\in\C| |z-1|=1}, F=\set{z\in\C|\Re(z)=\frac{1}{2}}$$
\begin{questions}
\item Déterminer la nature géométrique de $E$ et $F$.
\item Montrer que $f(E\setminus\set{0})\subset F$.
\item Montrer que $f$ induit une bijection de $E$ sur $F$.
\end{questions}
\end{questions}
\end{exercise}

\begin{exercise}
Soient $E$ un ensemble non vide et $f:E\to E$. 
Pour tous $A,B\in P(E)$, on note $A\Delta B=(A\cup B)\setminus(A\cap B)$.
\begin{questions}
\item Montrer que $\forall A,B\in P(E), f^{-1}(A\Delta B)=f^{-1}(A)\Delta f^{-1}(B)$.
\item Montrer que $\forall A,B\in P(E), f(A)\Delta f(B)\subset f(A\Delta B)$. 
A quelle condition, sur $f$, a-t-on égalité?
\end{questions}
\end{exercise}

\begin{exercise}
Soit $f:\R\to \R$. On définit une relation $\le_f$ sur $\R$ par:
$$\forall x,y\in\R, x\le_f y \iff f(y)-f(x)\ge|y-x|$$
\begin{questions}
\item Montrer que $\le_f$ est une relation d'ordre.
\item Montrer que $\le_f$ est totale ssi:
$$\forall x,y\in\R, |f(y)-f(x)|\ge|y-x|$$
\item A quoi la relation $\le_{Id_{\R}}$ est-elle égale?
\end{questions}
\end{exercise}






\end{document}